\begin{landscape}
\captionsetup{hypcap=false}
\section*{Appendix}
\begin{center}
\scriptsize\setstretch{1.15}\setlength{\tabcolsep}{3pt}
\captionof{table}{Expanding Window Results: Bitcoin (BTC)}\label{tab:expanding_btc}
\begin{minipage}[t]{0.49\linewidth}\centering
\textit{Panel A: 1-Hour RV}\par\smallskip
\begin{adjustbox}{max width=\linewidth}\begin{tabular}{@{}llrrrr@{}}
\toprule
Window & Model & $R^2$ & DA & QLIKE & MSE \\
\midrule
\multirow{5}{*}{2023 Jan--Jun} & HAR & -0.253 & 0.500 & 1.207 & 2.940 \\
 & EGARCH & -0.001 & 0.615 & 1.038 & 2.349 \\
 & EGARCH+NED & 0.280 & 0.675 & 1.528 & 1.689 \\
 & TS-only & 0.378 & 0.717 & 1.323 & 1.460 \\
 & NED & 0.380 & 0.716 & 1.246$^{***}$\,(-5.9\%) & 1.455$^{**}$\,(-0.3\%) \\
\midrule
\multirow{5}{*}{2023 Jul--Dec} & HAR & -0.584 & 0.500 & 1.169 & 3.652 \\
 & EGARCH & 0.078 & 0.648 & 0.956 & 2.125 \\
 & EGARCH+NED & 0.381 & 0.731 & 1.423 & 1.428 \\
 & TS-only & 0.451 & 0.758 & 1.372 & 1.267 \\
 & NED & 0.454 & 0.762 & 1.281$^{*}$\,(-6.6\%) & 1.259$^{**}$\,(-0.6\%) \\
\midrule
\multirow{5}{*}{2024 Jan--Jun} & HAR & 0.043 & 0.500 & 1.188 & 1.854 \\
 & EGARCH & 0.088 & 0.641 & 0.724 & 1.766 \\
 & EGARCH+NED & 0.359 & 0.705 & 1.059 & 1.241 \\
 & TS-only & 0.438 & 0.735 & 0.890 & 1.089 \\
 & NED & 0.449 & 0.737 & 0.772$^{***}$\,(-13.2\%) & 1.068$^{***}$\,(-1.9\%) \\
\midrule
\multirow{5}{*}{2024 Jul--Dec} & HAR & 0.060 & 0.550 & 1.171 & 1.655 \\
 & EGARCH & -0.020 & 0.609 & 0.677 & 1.796 \\
 & EGARCH+NED & 0.268 & 0.673 & 0.964 & 1.290 \\
 & TS-only & 0.352 & 0.704 & 0.833 & 1.141 \\
 & NED & 0.361 & 0.703 & 0.776$^{***}$\,(-6.8\%) & 1.125$^{***}$\,(-1.4\%) \\
\midrule
\multirow{5}{*}{2025 Jan--Jun} & HAR & 0.035 & 0.500 & 1.143 & 1.835 \\
 & EGARCH & 0.110 & 0.647 & 0.674 & 1.694 \\
 & EGARCH+NED & 0.382 & 0.726 & 0.941 & 1.175 \\
 & TS-only & 0.462 & 0.750 & 0.749 & 1.024 \\
 & NED & 0.468 & 0.751 & 0.693$^{***}$\,(-7.5\%) & 1.011$^{***}$\,(-1.3\%) \\
\midrule
\multirow{5}{*}{2025 Jul--Dec} & HAR & -0.136 & 0.500 & 0.939 & 2.094 \\
 & EGARCH & 0.045 & 0.643 & 0.688 & 1.760 \\
 & EGARCH+NED & 0.346 & 0.706 & 1.004 & 1.206 \\
 & TS-only & 0.421 & 0.735 & 0.784 & 1.067 \\
 & NED & 0.425 & 0.735 & 0.743$^{***}$\,(-5.3\%) & 1.060$^{***}$\,(-0.7\%) \\
\bottomrule
\end{tabular}
\end{adjustbox}
\end{minipage}\hfill
\begin{minipage}[t]{0.49\linewidth}\centering
\textit{Panel B: 4-Hour RV}\par\smallskip
\begin{adjustbox}{max width=\linewidth}\begin{tabular}{@{}llrrrr@{}}
\toprule
Window & Model & $R^2$ & DA & QLIKE & MSE \\
\midrule
\multirow{5}{*}{2023 Jan--Jun} & HAR & -0.338 & 0.500 & 0.769 & 2.175 \\
 & EGARCH & 0.137 & 0.639 & 0.783 & 1.403 \\
 & EGARCH+NED & 0.244 & 0.654 & 1.041 & 1.229 \\
 & TS-only & 0.370 & 0.684 & 0.806 & 1.024 \\
 & NED & 0.382 & 0.689 & 0.737$^{***}$\,(-8.5\%) & 1.005$^{***}$\,(-1.8\%) \\
\midrule
\multirow{5}{*}{2023 Jul--Dec} & HAR & -0.868 & 0.500 & 0.915 & 3.118 \\
 & EGARCH & 0.220 & 0.677 & 0.735 & 1.302 \\
 & EGARCH+NED & 0.367 & 0.710 & 0.949 & 1.057 \\
 & TS-only & 0.439 & 0.741 & 0.940 & 0.936 \\
 & NED & 0.457 & 0.747 & 0.867$^{***}$\,(-7.8\%) & 0.906$^{***}$\,(-3.2\%) \\
\midrule
\multirow{5}{*}{2024 Jan--Jun} & HAR & 0.013 & 0.500 & 0.685 & 1.340 \\
 & EGARCH & 0.232 & 0.668 & 0.506 & 1.043 \\
 & EGARCH+NED & 0.363 & 0.706 & 0.671 & 0.865 \\
 & TS-only & 0.438 & 0.732 & 0.578 & 0.763 \\
 & NED & 0.469 & 0.735 & 0.436$^{***}$\,(-24.6\%) & 0.721$^{***}$\,(-5.5\%) \\
\midrule
\multirow{5}{*}{2024 Jul--Dec} & HAR & 0.069 & 0.500 & 0.613 & 1.034 \\
 & EGARCH & 0.121 & 0.629 & 0.453 & 0.976 \\
 & EGARCH+NED & 0.278 & 0.653 & 0.556 & 0.803 \\
 & TS-only & 0.379 & 0.682 & 0.479 & 0.690 \\
 & NED & 0.402 & 0.688 & 0.428$^{***}$\,(-10.5\%) & 0.664$^{***}$\,(-3.7\%) \\
\midrule
\multirow{5}{*}{2025 Jan--Jun} & HAR & 0.000 & 0.500 & 0.685 & 1.356 \\
 & EGARCH & 0.246 & 0.666 & 0.502 & 1.022 \\
 & EGARCH+NED & 0.399 & 0.724 & 0.566 & 0.816 \\
 & TS-only & 0.470 & 0.732 & 0.511 & 0.719 \\
 & NED & 0.490 & 0.741 & 0.474$^{***}$\,(-7.4\%) & 0.692$^{***}$\,(-3.8\%) \\
\midrule
\multirow{5}{*}{2025 Jul--Dec} & HAR & -0.269 & 0.500 & 0.598 & 1.588 \\
 & EGARCH & 0.188 & 0.683 & 0.478 & 1.017 \\
 & EGARCH+NED & 0.353 & 0.721 & 0.557 & 0.810 \\
 & TS-only & 0.420 & 0.730 & 0.515 & 0.727 \\
 & NED & 0.426 & 0.738 & 0.473$^{***}$\,(-8.2\%) & 0.719\,(-1.0\%) \\
\bottomrule
\end{tabular}
\end{adjustbox}
\end{minipage}
\par\medskip
\begin{minipage}{\linewidth}\footnotesize
\textit{Notes:} Definitions follow Table~\ref{tab:forecast_summary}.
Parentheses on the NED row give percentage loss changes relative to
TS-only; negative means improvement. $^*p<0.10$, $^{**}p<0.05$,
$^{***}p<0.01$ from the corrected DM test within each window.
\end{minipage}
\end{center}
\newpage
\begin{center}
\scriptsize\setstretch{1.15}\setlength{\tabcolsep}{3pt}
\captionof{table}{Expanding Window Results: Ethereum (ETH)}\label{tab:expanding_eth}
\begin{minipage}[t]{0.49\linewidth}\centering
\textit{Panel A: 1-Hour RV}\par\smallskip
\begin{adjustbox}{max width=\linewidth}\begin{tabular}{@{}llrrrr@{}}
\toprule
Window & Model & $R^2$ & DA & QLIKE & MSE \\
\midrule
\multirow{5}{*}{2023 Jan--Jun} & HAR & -0.589 & 0.500 & 1.099 & 3.131 \\
 & EGARCH & -0.156 & 0.585 & 1.023 & 2.278 \\
 & EGARCH+NED & 0.170 & 0.659 & 1.476 & 1.635 \\
 & TS-only & 0.269 & 0.689 & 1.238 & 1.439 \\
 & NED & 0.262 & 0.684 & 1.191$^{***}$\,(-3.8\%) & 1.453$^{***}$\,(+1.0\%) \\
\midrule
\multirow{5}{*}{2023 Jul--Dec} & HAR & -0.966 & 0.500 & 1.209 & 4.089 \\
 & EGARCH & -0.024 & 0.619 & 0.967 & 2.130 \\
 & EGARCH+NED & 0.304 & 0.708 & 1.458 & 1.446 \\
 & TS-only & 0.375 & 0.735 & 1.398 & 1.300 \\
 & NED & 0.378 & 0.736 & 1.355$^{***}$\,(-3.1\%) & 1.294$^{**}$\,(-0.4\%) \\
\midrule
\multirow{5}{*}{2024 Jan--Jun} & HAR & -0.030 & 0.500 & 1.018 & 1.810 \\
 & EGARCH & -0.006 & 0.607 & 0.727 & 1.768 \\
 & EGARCH+NED & 0.279 & 0.679 & 1.121 & 1.268 \\
 & TS-only & 0.349 & 0.706 & 0.963 & 1.144 \\
 & NED & 0.350 & 0.706 & 0.949$^{***}$\,(-1.4\%) & 1.144\,(-0.048\%) \\
\midrule
\multirow{5}{*}{2024 Jul--Dec} & HAR & 0.074 & 0.500 & 1.103 & 1.523 \\
 & EGARCH & -0.041 & 0.605 & 0.686 & 1.713 \\
 & EGARCH+NED & 0.234 & 0.667 & 1.014 & 1.261 \\
 & TS-only & 0.320 & 0.703 & 0.909 & 1.119 \\
 & NED & 0.325 & 0.704 & 0.873$^{***}$\,(-3.9\%) & 1.111$^{***}$\,(-0.7\%) \\
\midrule
\multirow{5}{*}{2025 Jan--Jun} & HAR & 0.105 & 0.629 & 1.482 & 1.445 \\
 & EGARCH & 0.040 & 0.634 & 0.659 & 1.551 \\
 & EGARCH+NED & 0.278 & 0.695 & 1.010 & 1.166 \\
 & TS-only & 0.357 & 0.722 & 0.829 & 1.039 \\
 & NED & 0.356 & 0.722 & 0.856$^{***}$\,(+3.3\%) & 1.039\,(+0.1\%) \\
\midrule
\multirow{5}{*}{2025 Jul--Dec} & HAR & 0.064 & 0.624 & 1.207 & 1.541 \\
 & EGARCH & -0.004 & 0.616 & 0.738 & 1.653 \\
 & EGARCH+NED & 0.252 & 0.674 & 1.136 & 1.232 \\
 & TS-only & 0.338 & 0.706 & 0.893 & 1.089 \\
 & NED & 0.340 & 0.707 & 0.921$^{***}$\,(+3.2\%) & 1.087$^{*}$\,(-0.2\%) \\
\bottomrule
\end{tabular}
\end{adjustbox}
\end{minipage}\hfill
\begin{minipage}[t]{0.49\linewidth}\centering
\textit{Panel B: 4-Hour RV}\par\smallskip
\begin{adjustbox}{max width=\linewidth}\begin{tabular}{@{}llrrrr@{}}
\toprule
Window & Model & $R^2$ & DA & QLIKE & MSE \\
\midrule
\multirow{5}{*}{2023 Jan--Jun} & HAR & -0.791 & 0.500 & 0.758 & 2.378 \\
 & EGARCH & -0.012 & 0.615 & 0.741 & 1.343 \\
 & EGARCH+NED & 0.117 & 0.628 & 0.976 & 1.173 \\
 & TS-only & 0.236 & 0.659 & 0.684 & 1.014 \\
 & NED & 0.233 & 0.659 & 0.662$^{***}$\,(-3.2\%) & 1.019$^{**}$\,(+0.5\%) \\
\midrule
\multirow{5}{*}{2023 Jul--Dec} & HAR & -1.401 & 0.500 & 0.994 & 3.474 \\
 & EGARCH & 0.154 & 0.675 & 0.693 & 1.224 \\
 & EGARCH+NED & 0.321 & 0.712 & 0.892 & 0.983 \\
 & TS-only & 0.362 & 0.743 & 0.913 & 0.923 \\
 & NED & 0.360 & 0.743 & 0.860$^{***}$\,(-5.9\%) & 0.927\,(+0.3\%) \\
\midrule
\multirow{5}{*}{2024 Jan--Jun} & HAR & -0.086 & 0.500 & 0.590 & 1.220 \\
 & EGARCH & 0.182 & 0.662 & 0.473 & 0.919 \\
 & EGARCH+NED & 0.313 & 0.693 & 0.645 & 0.771 \\
 & TS-only & 0.383 & 0.711 & 0.582 & 0.692 \\
 & NED & 0.380 & 0.714 & 0.570$^{***}$\,(-2.2\%) & 0.696\,(+0.6\%) \\
\midrule
\multirow{5}{*}{2024 Jul--Dec} & HAR & 0.095 & 0.500 & 0.572 & 0.931 \\
 & EGARCH & 0.136 & 0.645 & 0.424 & 0.889 \\
 & EGARCH+NED & 0.253 & 0.667 & 0.541 & 0.768 \\
 & TS-only & 0.352 & 0.699 & 0.520 & 0.666 \\
 & NED & 0.369 & 0.703 & 0.482$^{***}$\,(-7.3\%) & 0.649$^{***}$\,(-2.6\%) \\
\midrule
\multirow{5}{*}{2025 Jan--Jun} & HAR & 0.143 & 0.663 & 0.810 & 0.936 \\
 & EGARCH & 0.236 & 0.683 & 0.443 & 0.834 \\
 & EGARCH+NED & 0.323 & 0.709 & 0.586 & 0.739 \\
 & TS-only & 0.414 & 0.732 & 0.514 & 0.640 \\
 & NED & 0.404 & 0.734 & 0.539$^{***}$\,(+4.7\%) & 0.650$^{**}$\,(+1.6\%) \\
\midrule
\multirow{5}{*}{2025 Jul--Dec} & HAR & 0.074 & 0.621 & 0.640 & 0.981 \\
 & EGARCH & 0.109 & 0.645 & 0.520 & 0.944 \\
 & EGARCH+NED & 0.212 & 0.674 & 0.679 & 0.834 \\
 & TS-only & 0.354 & 0.705 & 0.523 & 0.685 \\
 & NED & 0.352 & 0.705 & 0.548$^{***}$\,(+4.8\%) & 0.686\,(+0.2\%) \\
\bottomrule
\end{tabular}
\end{adjustbox}
\end{minipage}
\par\medskip
\begin{minipage}{\linewidth}\footnotesize
\textit{Notes:} Metric definitions, percentage-change conventions,
and significance levels are identical to Table~\ref{tab:expanding_btc}.
\end{minipage}
\end{center}
\end{landscape}

\subsection*{Sensitivity to lookback horizon}
We repeat the BTC four-hour exponential-decay experiment with lookbacks
of six, twelve, and twenty-four hours, keeping the P1 baseline fixed
and refitting P2 at each horizon. All comparisons use the same six
test periods and the same training objective, learning rate, batch
size, and stopping rule. This separates a change in the available
news history from a change in the time-series model.

\begin{table}[htbp]
\centering\small
\caption{Sensitivity to the news lookback (4-Hour RV, BTC)}\label{tab:lookback}
\begin{tabular}{lrrrr}
\toprule
Lookback & $R^2$ & DA & QLIKE & MSE \\
\midrule
6 hours & 0.437 & 0.724 & 0.559 & 0.786 \\
12 hours & 0.434 & 0.723 & 0.571 & 0.790 \\
24 hours & 0.436 & 0.723 & 0.566 & 0.786 \\
\bottomrule
\end{tabular}

\tabnote{\textit{Notes:} Arithmetic means across six windows. The
twelve-hour specification is the refitted exponential reference
reported in Table~\ref{tab:decay_ablation}. These are separate P2
fits, not test-time truncations of a fixed twelve-hour model.}
\end{table}

The forecast improvement persists across all three lookbacks in
Table~\ref{tab:lookback}. Six hours produces the lowest mean QLIKE,
0.55898, and log-MSE, 0.78577. At twelve hours the losses are
0.57071 and 0.78967, while extending the window to twenty-four
hours gives 0.56564 and 0.78618. The six-hour and twenty-four-hour
log-MSE values are particularly close. Thus the main result is not
confined to a twelve-hour news history, and a longer lookback does
not produce a monotonic improvement. These separately fitted
single-seed comparisons do not identify an optimal lookback beyond
the evaluated sample; choosing a new default would require further
validation.

